% §9 Conclusion
% Status: Written 2026-03-26 — full rewrite with RQ answers and empirical evidence.

\section{Conclusion}
\label{sec:conclusion}

Abstraction selection in GPU-accelerated HPC software has long been guided
by heuristic judgment, developer experience, and trial-and-error. This paper
presented a measurement-driven framework for \emph{selective abstraction}
that replaces this opaque process with evidence-based prediction. Through
a systematic evaluation of four abstraction layers (Kokkos, RAJA, SYCL,
Julia) against native CUDA and HIP baselines across seven representative
workloads and two GPU architectures (NVIDIA RTX~5060 and AMD MI300X), we
measured over 190 configurations and constructed a quantified decision
framework grounded in empirical observation.

% ------------------------------------------------------------------
\subsection*{Research Question Answers}
% ------------------------------------------------------------------

\paragraph{RQ1: When does abstraction help?}
Abstraction delivers near-native performance when kernel execution time
substantially exceeds per-launch overhead and the abstraction's API can
express the optimal algorithm. Specifically, when $T_k > 10 \times
\text{overhead}$ and no API expressivity gap exists (R5), all four
abstractions achieve PPC $\geq 0.90$ (Section~\ref{sec:framework}). At
large problem sizes with memory-bound workloads, all abstractions achieve
PPC $\geq 0.91$ on both platforms (E1, Section~\ref{subsec:e1}). RAJA and
Kokkos can even \emph{exceed} native performance on wide-level workloads
through synchronisation pipelining advantages (RAJA: PPC $= 1.299$ on E5
random/large, Section~\ref{subsec:e5}). Across the full experimental
matrix, RAJA achieves the best cross-platform portability
($\varphi = 0.954$), followed by Kokkos ($\varphi = 0.886$).

\paragraph{RQ2: When does abstraction fail?}
Abstraction fails in structurally predictable ways. We identified eight
patterns (P001--P008), of which five are validated with quantitative
evidence (Section~\ref{sec:taxonomy}). Three patterns --- launch overhead
dominance (P001), API expressivity gap (P004), and level-set dispatch
amplification (P008) --- account for approximately 80\% of all
configurations with PPC $< 0.60$. These failures are not random: they
recur whenever the same combination of workload structure and overhead
profile is present, independent of framework version. Julia is the most
architecture-sensitive abstraction: its column-major layout advantage
produces PPC $= 1.25$ on NVIDIA DGEMM (P005) but reverses to PPC $= 0.31$
on AMD, where the native baseline is $4\times$ stronger. SYCL's
performance variance is explained almost entirely by workload structure:
single-kernel workloads achieve 0.92--0.99 efficiency, while multi-kernel
level-set workloads collapse to 0.30--0.49 through HSA dispatch
amplification.

\paragraph{RQ3: How can we quantify the trade-off?}
Yes. Five calibrated decision rules (R1--R5), derived from five measurable
workload characteristics ($T_k$, $L$, $\sigma$, $AI$,
$\text{API}_{\text{ok}}$), predict abstraction PPC with 85\%
tier-correct accuracy on 39 held-out configurations from E6 (BFS) and
E7 (N-Body) that were not used during rule calibration
(Section~\ref{sec:validation}). The strongest validation is R2's exact
prediction of SYCL/SpTRSV performance on AMD MI300X (predicted PPC
$= 0.490$, observed $= 0.490$). The framework is implemented as a
command-line tool (\texttt{abstraction-advisor}) that ingests workload
profiles and produces ranked abstraction recommendations with pattern
warnings.

% ------------------------------------------------------------------
\subsection*{Contributions}
% ------------------------------------------------------------------

This paper makes five contributions:

\begin{enumerate}
    \item An \textbf{empirical taxonomy} of eight abstraction failure and
    success patterns (P001--P008), five validated with structural recurrence
    across experiments (Section~\ref{sec:taxonomy}).

    \item A \textbf{measurement methodology} comprising a four-level
    hierarchy (PPC $\rightarrow$ overhead attribution $\rightarrow$ root
    cause analysis $\rightarrow$ roofline characterisation) applied across
    190+ configurations (Section~\ref{sec:methodology}).

    \item A \textbf{quantified decision framework} of five calibrated rules
    with 85\% tier-correct predictive accuracy on held-out workloads
    (Section~\ref{sec:framework}).

    \item A \textbf{performance portability dataset} spanning 4 abstraction
    layers, 2 GPU architectures, and 7 workloads, with full reproducibility
    artifacts.

    \item A \textbf{decision support tool}
    (\texttt{abstraction-advisor}) that operationalises the framework as a
    command-line advisor for abstraction selection.
\end{enumerate}

% ------------------------------------------------------------------
\subsection*{Limitations}
% ------------------------------------------------------------------

This work is subject to several limitations that bound the scope of its
claims:

\begin{itemize}
    \item \textbf{Two hardware platforms.} The experimental matrix covers
    NVIDIA RTX~5060 (Blackwell, consumer-grade) and AMD MI300X (CDNA3, HPC).
    Intel GPU Max, NVIDIA H100/GH200, and AMD CDNA4 are not evaluated.
    The RTX~5060's consumer-grade classification limits the
    representativeness of the NVIDIA results for HPC-class deployments,
    though the framework's rule structure is expected to transfer with
    recalibrated overhead thresholds.

    \item \textbf{Seven kernels.} The workload set covers regular,
    semi-irregular, and irregular patterns but does not include FFT, AMR,
    particle-in-cell, or deep learning workloads. Patterns not represented
    in the training set cannot be predicted.

    \item \textbf{Single-GPU scope.} Multi-GPU execution and
    distributed-memory parallelism introduce overhead dimensions
    (inter-GPU communication, distributed load balancing) outside the
    current model.

    \item \textbf{Temporal snapshot.} Per-launch overhead thresholds are
    tied to specific compiler, framework, and driver versions. The
    calibration protocol is reproducible, but threshold values will require
    recalibration as software stacks evolve.

    \item \textbf{Unconfirmed patterns.} P002 (Compiler Backend Failure)
    and P003 (Memory Layout Mismatch) could not be confirmed due to
    hardware profiling constraints (\texttt{ncu} blocked on RTX~5060).
\end{itemize}

% ------------------------------------------------------------------
\subsection*{Future Work}
% ------------------------------------------------------------------

Several directions extend this work:

\begin{itemize}
    \item \textbf{Hardware expansion.} Recalibrating overhead thresholds
    for Intel GPU Max~1550, NVIDIA H100/GH200, and AMD CDNA4 would test
    whether the rule structure transfers across hardware generations with
    only threshold updates.

    \item \textbf{Workload expansion.} Adding FFT, AMR, particle-in-cell,
    and deep learning inference kernels would test whether the eight-pattern
    taxonomy is complete or whether additional failure modes emerge.

    \item \textbf{Multi-GPU and distributed execution.} Extending the
    framework to model MPI communication overhead, inter-GPU data transfer,
    and distributed load balancing would address the single-node limitation.

    \item \textbf{Retrospective ECP validation.} Applying the framework
    retrospectively to production ECP codes (ExaSMR, QMCPACK, LAMMPS) would
    test whether it predicts the abstraction choices made by expert
    developers in practice.

    \item \textbf{Living taxonomy.} Establishing a community-maintained
    pattern database with standardised contribution protocols would enable
    continuous threshold recalibration and pattern discovery as new
    frameworks and architectures emerge.

    \item \textbf{Automated root-cause attribution.} Integrating PTX/GCN
    inspection and hardware counter analysis into the decision tool would
    enable automated diagnosis of patterns P002 and P003, which currently
    require manual profiling.
\end{itemize}

\medskip
\noindent Selective abstraction --- choosing the right level of abstraction
for each workload, architecture, and performance target --- is not an art
but an engineering discipline amenable to measurement, classification, and
prediction. This paper provides the methodology and evidence base to
practise it.
